% Solutions for appendix B-blowup
\providecommand{\Vb}{\mathcal{V}_{\mathrm{b}}}
\providecommand{\bT}{{}^{\mathrm{b}}T}
\providecommand{\relRC}[2]{{}^{#1}\overline{#2}}

\section[Solutions: real blow-up]{Solutions for \cref{ch:blowup}: Real blow-up and manifolds with corners}\label{solutions:blowup}

%----------------------------------------------------------------------
\subsection*{\Cref{blowup:ex:projective}}
Idea: everything follows from the chain rule for the substitution $z_1=s$, $z_k=st_k$.
Throughout $n\ge2$, as in the exercise: the coordinate $t_2$ and the rotation
$z_1\pa_{z_2}-z_2\pa_{z_1}$ exist only then, and for $n=1$ the chart has no $t$ variables and
$\pa_{z_1}$ lifts to the smooth field $\pa_s$ (compare \cref{blowup:lem:vf-point}).

(a) On $\{z_1>0\}$ the map $(s,t)\mapsto z=(s,st_2,\dots,st_n)$ is a diffeomorphism onto
$\{z_1>0\}$, with inverse $s=z_1$, $t_k=z_k/z_1$. A vector field is determined by its
action on the coordinate functions $s$, $t_k$. Since $\pa_{z_k}s=0$, $\pa_{z_k}t_r=\delta_{kr}/z_1$
for $k\ge2$, and $\pa_{z_1}s=1$, $\pa_{z_1}t_r=-z_r/z_1^2=-t_r/s$, we get
\begin{equation*}
  \pa_{z_k}=s^{-1}\pa_{t_k}\ (k\ge2),\qquad\pa_{z_1}=\pa_s-s^{-1}\sum_{r\ge2}t_r\pa_{t_r}.
\end{equation*}
Multiplying by $z_l=st_l$ ($l\ge2$) or $z_1=s$ gives the four cases of
\eqref{blowup:eq:proj-lifts} with $j=1$: $z_l\pa_{z_m}=t_l\pa_{t_m}$ and $z_1\pa_{z_m}=\pa_{t_m}$
for $m\ge2$; $z_l\pa_{z_1}=t_l(s\pa_s-\sum_rt_r\pa_{t_r})$ and
$z_1\pa_{z_1}=s\pa_s-\sum_rt_r\pa_{t_r}$. All are smooth up to $s=0$, which is the front face
by \cref{blowup:lem:projective}.

(b) Summing, $\sum_lz_l\pa_{z_l}=(s\pa_s-\sum_{r\ge2}t_r\pa_{t_r})+\sum_{l\ge2}t_l\pa_{t_l}=s\pa_s$.
For the rotation,
\begin{equation*}
  z_1\pa_{z_2}-z_2\pa_{z_1}=\pa_{t_2}-t_2\Big(s\pa_s-\sum_{r\ge2}t_r\pa_{t_r}\Big)
  =(1+t_2^2)\pa_{t_2}-t_2s\pa_s+t_2\sum_{r\ge3}t_r\pa_{t_r}.
\end{equation*}
It is tangent to $\{s=0\}$ but does not vanish there.

(c) By (a), $\pa_{z_1}=s^{-1}V$ with $V=s\pa_s-\sum_{r\ge2}t_r\pa_{t_r}$ smooth. If $\pa_{z_1}$
had a smooth extension $\tilde V$ near a point $(0,t^0)$ of the front face, then $\tilde Vt_2$
would be a smooth, hence locally bounded, function. But for $s>0$, $\tilde Vt_2=\pa_{z_1}t_2=-t_2/s$,
which is unbounded as $s\downarrow0$ along the curve $t_2=t^0_2+s^{1/2}$, $t_k=t^0_k$ ($k\ge3$).
So $\pa_{z_1}$ has no smooth lift near any point of the front face in this chart.

(d) The Jacobian matrix of $(s,t)\mapsto z$ is lower triangular, with first row $(1,0,\dots,0)$
and $dz_k=t_k\,ds+s\,dt_k$, so its determinant is $s^{n-1}$. Hence
$\beta^*\abs{dz}=s^{n-1}\abs{ds\,dt}$, in agreement with \cref{blowup:prop:density}.

%----------------------------------------------------------------------
\subsection*{\Cref{blowup:ex:cone}}
Idea: the cone is invariant under dilations, so in polar coordinates it is a product.

Use polar coordinates $(r,\omega)$, $\omega=(\omega_t,\omega_x,\omega_y)\in\Sphere^2$, on
$[\R^3;\{0\}]=[0,\infty)\times\Sphere^2$. For $r>0$, $r\omega\in C$ if and only if
$\omega_t^2=\omega_x^2+\omega_y^2$ and $\omega_t\ge0$; with $\abs\omega=1$ this means
$\omega_t=1/\sqrt2$ and $(\omega_x,\omega_y)=\theta/\sqrt2$, $\theta\in\Sphere^1$. Let
$C_1=\{\omega\in\Sphere^2;\ \omega_t=1/\sqrt2\}$. Then $\beta^{-1}(C\setminus0)=(0,\infty)\times C_1$,
whose closure is $\tilde C=[0,\infty)\times C_1$.

$C_1$ is a smooth circle in $\Sphere^2$, since the restriction of $\omega_t$ to $\Sphere^2$
has non-zero differential where $\abs{\omega_t}<1$. Near a point of $C_1$ choose an angular
coordinate $\vartheta$ on $\Sphere^2$ with $(\omega_t-1/\sqrt2,\vartheta)$ local coordinates; then
$(r,\omega_t-1/\sqrt2,\vartheta)$ are coordinates on $[0,\infty)\times\Sphere^2$ in which
$\tilde C=\{\omega_t-1/\sqrt2=0\}$ with $\omega_t-1/\sqrt2$ an interior coordinate. So $\tilde C$
is an interior p-submanifold, diffeomorphic to $[0,\infty)\times\Sphere^1$ by
$(r,\theta)\mapsto(r,(1,\theta)/\sqrt2)$. It meets the front face $\{r=0\}$ in $\{0\}\times C_1$,
transversally since $dr\ne0$ on $\tilde C$. Finally $\beta(r,(1,\theta)/\sqrt2)=r(1,\theta)/\sqrt2$.
Thus the singular space $C$ is resolved to a manifold with boundary, the boundary being
the circle of directions in which $C$ leaves the vertex.

%----------------------------------------------------------------------
\subsection*{\Cref{blowup:ex:morse}}
Idea: $u-u(0)$ is homogeneous of degree two to leading order, so it is $r^2$ times a
function whose zero set on the front face is the cone of null directions of the Hessian.

Substituting $z=r\omega$ gives $\beta^*u=u(0)+r^2\omega^TH(r\omega)\omega=u(0)+r^2U$, with
$U(r,\omega)=\omega^TH(r\omega)\omega$ smooth on $[0,\infty)\times\Sphere^{n-1}$. At $r=0$,
$U(0,\omega)=q(\omega)=\omega^TH(0)\omega$. The differential of $q$ along the sphere at $\omega$ is the
tangential part of $\nabla q(\omega)=2H(0)\omega$, namely $2H(0)\omega-2(\omega\cdot H(0)\omega)\omega$. At a
point where $q(\omega)=0$ this equals $2H(0)\omega$, which is non-zero because $H(0)$ is invertible.
Thus $d_\omega U\ne0$ at every point of the compact set $\{r=0,\ U=0\}$, and hence, by continuity,
$dU\ne0$ on a neighbourhood $\mathcal N$ of it.

The full inverse image of the level set is
$\beta^{-1}(\{u=u(0)\})=\{\beta^*u=u(0)\}=\{r^2U=0\}=\{r=0\}\cup\{U=0\}$, the union of the whole
front face and $\{U=0\}$. In $\mathcal N$,
$\{U=0\}$ is a smooth hypersurface, and it meets $\{r=0\}$ transversally: at a common point
$dU$ has a non-zero component along the sphere, while $dr$ has none, so $dr$ and $dU$ are
independent. (If $H(0)$ is definite, $q\ne0$ on the sphere and near the front face the lifted
level set is just the front face, reflecting that $0$ is an isolated point of the level set.)
In other words, after one blow-up the full inverse image of the singular level set has normal
crossings.

The lift in the sense of \cref{blowup:def:lift} (the strict transform) is different: it is the
closure of $\beta^{-1}(\{u=u(0)\}\setminus0)=\{r>0,\ U=0\}$. Near the front face this is
$\{U=0\}\cap\mathcal N$, because $\{U=0\}$ is a hypersurface transversal to $\{r=0\}$, so its
part with $r>0$ is dense in it; it does not contain the rest of the front face. It is a
p-submanifold (with $U$ as a normal, interior coordinate) meeting the front face transversally in
$\{r=0,\ q(\omega)=0\}$, the set of null directions of $H(0)$. For example, for
$u=z_1^2-z_2^2$ on $\R^2$ the strict transform is the union of the closures of the lifts of the
four rays $z_2=\pm z_1$, i.e.\ $[0,\infty)\times\{\omega;\ \omega_1=\pm\omega_2\}$, four half-lines meeting
the front face circle in four points, while the full inverse image also contains the whole
front face.

%----------------------------------------------------------------------
\subsection*{\Cref{blowup:ex:diag-not-p}}
Idea: a coordinate subspace meets a codimension-two corner in codimension two, but the
diagonal meets the corner of $X^2$ in codimension one.

Let $p\in\pa X$ and $q=(p,p)$, and let $(x,y)$, $(x',y')$ be boundary coordinates near $p$ on the
two factors. The boundary hypersurfaces of $X^2$ through $q$ are $\{x=0\}$ and $\{x'=0\}$, and
$q$ has codimension two. Near $q$,
\begin{equation*}
  \Diag\cap\{x=0\}=\Diag\cap\{x'=0\}=\Diag\cap\{x=x'=0\}=\{(p',p');\ p'\in\pa X\},
\end{equation*}
a manifold of dimension $n-1$ through $q$.

Suppose $\Diag$ were a p-submanifold, with a chart $(z_1,\dots,z_{2n})$ centred at $q$ in which
$\Diag=\{z_A=0\}$, $\abs A=n$. By \cref{blowup:lem:bdf} the two boundary hypersurfaces through
$q$ are $\{z_1=0\}$ and $\{z_2=0\}$ for two boundary coordinates, say $\{x=0\}=\{z_1=0\}$ and
$\{x'=0\}=\{z_2=0\}$ near $q$. If $1\in A$, then $\Diag\subset\{x=0\}$ near $q$, which is false:
$(p',p')$ with $p'$ an interior point near $p$ lies in $\Diag$. Similarly $2\notin A$. But then
$\Diag\cap\{z_1=z_2=0\}=\{z_A=0,\ z_1=z_2=0\}$ is a coordinate subspace of dimension
$2n-n-2=n-2$ (empty if $n=1$, whereas it contains $q$), contradicting the computation above.
Hence $\Diag$ is not a p-submanifold. (\Cref{blowup:exa:b-double} shows that its lift to
$[X^2;(\pa X)^2]$ is one.)

%----------------------------------------------------------------------
\subsection*{\Cref{blowup:ex:noncommute}}
Idea: on $[\R^3;Y]$ the angle around $Y$ is a smooth function; on $[\R^3;Z;Y]$ the same
angle has different limits along two sequences converging to the same point.

\emph{Position.} In the coordinates $z$ themselves, $Y=\{z_A=0\}$ and $Z=\{z_B=0\}$ with
$A=\{1,2\}$, $B=\{2,3\}$, so they are joint. Since $A\cap B=\{2\}\ne\emptyset$ and neither of
$A$, $B$ contains the other, they are neither transversal nor nested. So
\cref{blowup:prop:lift} applies (both lifts are p-submanifolds), but
\cref{blowup:prop:commute} does not.

\emph{$f$ on $[\R^3;Y;Z]$.} On $[\R^3;Y]=[0,\infty)_\rho\times\Sphere^1_\phi\times\R_{z_3}$, with
$(z_1,z_2)=\rho(\cos\phi,\sin\phi)$, we have $f=\sin\phi$ for $\rho>0$, which extends smoothly.
Composing with the blow-down map $[\R^3;Y;Z]\to[\R^3;Y]$ gives a continuous (indeed smooth)
extension of $f$ to $[\R^3;Y;Z]$.

\emph{$f$ on $[\R^3;Z;Y]$.} On $[\R^3;Z]$, near directions with $z_3>0$, we use projective
coordinates $(z_1,\sigma,\kappa)$ with $\sigma=z_3\ge0$ and $\kappa=z_2/z_3$
(\cref{blowup:lem:projective}); the front face is $\{\sigma=0\}$. The lift of $Y$ is the closure
of $\{z_1=0,\ \kappa=0,\ \sigma>0\}$, which is $\{z_1=0,\ \kappa=0\}$ (\cref{blowup:prop:lift}).
Blowing it up we use polar coordinates $z_1=R\cos\alpha$, $\kappa=R\sin\alpha$, so
$(R,\alpha,\sigma)\in[0,\infty)\times\Sphere^1\times[0,\infty)$ are coordinates on an open set of
$[\R^3;Z;Y]$. Where $R,\sigma>0$ and $\sin\alpha\ne0$, $\cos\alpha\ne0$ we are in the interior and
$z_2=\sigma\kappa=\sigma R\sin\alpha$, so
\begin{equation*}
  f=\frac{\sigma R\sin\alpha}{\sqrt{R^2\cos^2\alpha+\sigma^2R^2\sin^2\alpha}}
  =\frac{\sigma\sin\alpha}{\sqrt{\cos^2\alpha+\sigma^2\sin^2\alpha}}.
\end{equation*}
Let $P$ be the point $(R,\alpha,\sigma)=(0,\pi/2,0)$. The points
$a_k=(1/k,\ \pi/2-1/k^2,\ 1/k)$ and $b_k=(1/k,\ \pi/2-1/k,\ 1/k^2)$ are interior points
converging to $P$. At $a_k$, $\cos\alpha\approx k^{-2}$ and $\sigma=k^{-1}$, so
$f(a_k)=\sigma\sin\alpha/\sqrt{\cos^2\alpha+\sigma^2\sin^2\alpha}\to1$ (the denominator is
$k^{-1}(1+O(k^{-2}))$). At $b_k$, $\cos\alpha\approx k^{-1}$ and $\sigma=k^{-2}$, so
$f(b_k)\approx k^{-2}/k^{-1}\to0$. Hence $f$ has no continuous extension to $[\R^3;Z;Y]$.

\emph{Conclusion.} If the identity of $\R^3\setminus(Y\cup Z)$ extended to a homeomorphism
$\Phi:[\R^3;Z;Y]\to[\R^3;Y;Z]$, then $\tilde f\circ\Phi$, with $\tilde f$ the continuous extension
on $[\R^3;Y;Z]$, would be a continuous extension of $f$ to $[\R^3;Z;Y]$, which does not exist.

%----------------------------------------------------------------------
\subsection*{\Cref{blowup:ex:corner-exponents}}
Idea: read off the exponents from \cref{blowup:prop:props}(ii), and the b-differential from
projective coordinates.

The boundary hypersurfaces of $[\R^{2,2};\{0\}]$ are the front face $\mathrm{ff}$ and the lifts
$H_1'$, $H_2'$ of $H_1=\{x_1=0\}$, $H_2=\{x_2=0\}$. Since $\{0\}\subset H_1\cap H_2$,
\cref{blowup:prop:props}(ii) gives $\beta^*x_i=\rho_{\mathrm{ff}}\rho_{H_i'}$, so the exponent matrix
(rows $\mathrm{ff},H_1',H_2'$, columns $H_1,H_2$) is
\begin{equation*}
  e=\begin{pmatrix}1&1\\1&0\\0&1\end{pmatrix}.
\end{equation*}
Explicitly, in polar coordinates $x=r(\cos\phi,\sin\phi)$, $\phi\in[0,\pi/2]$, we have
$\beta^*x_1=r\cos\phi$ and $\beta^*x_2=r\sin\phi$, where $r=\rho_{\mathrm{ff}}$; the lift $H_1'$ of
$\{x_1=0\}$ is the closure of $\{\phi=\pi/2,\ r>0\}$, with defining function $\cos\phi$, and
similarly $\sin\phi$ defines $H_2'=\{\phi=0\}$. The first row has two
non-zero entries: the front face is mapped into the corner, so $\beta$ is not b-normal.

For the b-submersion property, take the projective coordinates $s=x_1$, $t=x_2/x_1$ near
$\{\phi<\pi/2\}$. Then $\log\beta^*x_1=\log s$, $\log\beta^*x_2=\log s+\log t$, so by
\cref{blowup:lem:bdiff} $\beta_*(s\pa_s)=x_1\pa_{x_1}+x_2\pa_{x_2}$ and
$\beta_*(t\pa_t)=x_2\pa_{x_2}$, where $s\pa_s$, $t\pa_t$ is a basis of $\bT$ in this chart (both
$s,t\ge0$ are boundary coordinates). The images span $\bT\R^{2,2}$, and the other projective chart
$s=x_2$, $t=x_1/x_2$ is symmetric. Away from the front face $\beta$ is a diffeomorphism. So $\beta$
is a b-submersion (\cref{blowup:prop:bd-bsub}).

The composite $\pi_1\circ\beta$ is an interior b-map to $[0,\infty)$, whose only boundary
hypersurface is $\{0\}$; so each row of its exponent matrix has at most one entry and it is
b-normal. (Its exponents are $e(\mathrm{ff})=1$, $e(H_1')=1$, $e(H_2')=0$.)

In general, let $F$ be a boundary face of codimension $a$. Near a point of $F$ of
codimension $a$ in $X$, \cref{blowup:lem:bdf} shows that exactly $a$ boundary hypersurfaces pass
through it, namely $\{x_i=0\}$, $i\le a$, and they contain $F$; any hypersurface containing $F$
passes through this point. By \cref{blowup:prop:props}(ii), $e(\mathrm{ff},H)=1$ if $F\subset H$ and
$0$ otherwise. So the row of $\mathrm{ff}$ has exactly $a$ entries equal to $1$, and $\beta$ is b-normal
if and only if $a=1$, i.e.\ $[X;F]=X$.

%----------------------------------------------------------------------
\subsection*{\Cref{blowup:ex:half-disc}}
Idea: the fibres change type at the corners, where the restriction of $f$ to the curved face
is stationary.

\emph{$D_+$ is a manifold with corners.} Put $\rho_S=y$ and $\rho_A=1-x^2-y^2$. Then
$D_+=\{\rho_S\ge0,\ \rho_A\ge0\}$. At points where only one of them vanishes its differential is
non-zero ($dy$, resp.\ $-2(x\,dx+y\,dy)\ne0$ on the circle), so it is a boundary coordinate. At
$(\pm1,0)$ both vanish and $d\rho_S=dy$, $d\rho_A=\mp2\,dx$ are independent, so $(\rho_A,\rho_S)$ is a
chart onto a neighbourhood of $0$ in $\R^{2,2}$. Hence $D_+$ is a manifold with corners with two
boundary hypersurfaces, the segment $S=\{\rho_S=0\}$ and the arc $A=\{\rho_A=0\}$, embedded with
the defining functions $\rho_S$, $\rho_A$, and two corners $(\pm1,0)$.

\emph{$f$.} It is surjective onto $[-1,1]$ and $df=dx\ne0$ everywhere. Its fibre over
$x\in(-1,1)$ is the segment $\{x\}\times[0,\sqrt{1-x^2}]$, and over $\pm1$ the point $(\pm1,0)$.
If $f$ were a fibration, each point of $[-1,1]$ would have a neighbourhood over which the fibres
are all diffeomorphic; since $[-1,1]$ is connected, all fibres would be diffeomorphic. A point
is not diffeomorphic to a segment, so $f$ is not a fibration.

\emph{Which hypothesis fails.} Regard $f:D_+\to\R$. Its restriction to the face $A$ is
$x|_A$, whose differential vanishes at $(\pm1,0)$ (the arc is vertical there); and its restriction
to the zero-dimensional faces $\{(\pm1,0)\}$ is not a submersion onto $\R$. So the hypothesis of
\cref{blowup:prop:fibration} fails exactly at the corners. Consistently, $f(D_+)=[-1,1]$ is not
open in $\R$.

\emph{b-map?} The target $[-1,1]$ has boundary hypersurfaces $\{\pm1\}$ with defining functions
$1\mp x$. Now $f^*(1-x)=1-x$ vanishes on $D_+$ exactly at the point $(1,0)$. For an interior
b-map, $f^*(1-x)=a\rho_S^{e_S}\rho_A^{e_A}$ with $a>0$ would vanish either nowhere or on a union of
boundary hypersurfaces; and $f^*(1-x)\not\equiv0$. So $f$ is not a b-map.

%----------------------------------------------------------------------
\subsection*{\Cref{blowup:ex:octagon}}
Idea: count the faces created by each blow-up, and identify faces through the limits of curves
in $W$.

\emph{$\mathcal D_1$.} For $\dim U=\dim V=1$, $\RC W$ is a closed disc and $\Sphere U$, $\Sphere V$ are
pairs of antipodal points of the boundary circle, four points in all. Each is a p-submanifold
of normal type $(1,1)$ (one boundary and one interior normal variable), so blowing it up replaces
it by $\Sphere^{1,1}$, a closed half-circle, i.e.\ an interval, and cuts the circle there. The
resulting space has eight boundary hypersurfaces: four arcs of the circle (the components of the
lift of $\Sphere W$) and four front faces (two components each of $\mathrm{ff}_U$ and $\mathrm{ff}_V$),
alternating around the boundary, with eight corners. Topologically it is a disc with boundary
divided into eight intervals: an octagon.

\emph{$\mathcal D_2$.} $\RC V\times\RC U$ is a square with four sides
$(\Sphere V)\times\RC U$ (two sides) and $\RC V\times\Sphere U$ (two sides) and four corners
$\Sphere V\times\Sphere U$. Blowing up a corner (normal type $(2,0)$) inserts a quarter circle
$\Sphere^{2,0}$. Again we get eight sides alternating between the lifted sides of the square and
the four front faces.

\emph{Identification.} By \cref{blowup:prop:doubly} there is a diffeomorphism $\mathcal D_1\to\mathcal
D_2$ extending the identity of $W$; it maps boundary hypersurfaces onto boundary hypersurfaces,
and we can identify them by the limits of curves in $W$. The interior of a lifted side
$(\pm\Sphere V)\times\RC U$ of the square is reached by curves $(v,u)$ with $v\to\pm\infty$ and $u$
bounded; in $\RC W$ they converge to $\pm\Sphere V$, so in $\mathcal D_1$ they end on $\mathrm{ff}_V$.
Similarly the sides $\RC V\times(\pm\Sphere U)$ correspond to the components of $\mathrm{ff}_U$. The
interior of the front face over a corner of the square is reached by curves with $\abs u$,
$\abs v\to\infty$ and $\ang u/\ang v\to c\in(0,\infty)$; in $\RC W$ these converge to points of
$\Sphere W\setminus(\Sphere U\cup\Sphere V)$, so in $\mathcal D_1$ they end on the lift of $\Sphere W$. Thus:
\begin{equation*}
  \mathrm{ff}_V\leftrightarrow\text{lift of }\Sphere V\times\RC U,\qquad
  \mathrm{ff}_U\leftrightarrow\text{lift of }\RC V\times\Sphere U,\qquad
  \text{lift of }\Sphere W\leftrightarrow\mathrm{ff}(\Sphere V\times\Sphere U).
\end{equation*}
This is the picture of \cref{blowup:fig:octagon}.

\emph{Higher dimensions.} If $\dim U=k\ge2$ and $\dim V=l\ge2$, the spheres $\Sphere U$, $\Sphere
V$ and $\Sphere W$ are connected. By \cref{blowup:prop:props}(ii) the boundary hypersurfaces of
$\mathcal D_1$ are the lift of $\Sphere W$ and the front faces $\mathrm{ff}_U$, $\mathrm{ff}_V$. The front
face $\mathrm{ff}_V$ is a bundle over the connected manifold $\Sphere V$ with connected fibre
$\Sphere^{1,k}$ (a closed half-sphere), hence connected; likewise $\mathrm{ff}_U$. The lift of $\Sphere W$
is the closure of $\Sphere W\setminus(\Sphere U\cup\Sphere V)$, and the complement in the connected
manifold $\Sphere W$ of the two disjoint submanifolds $\Sphere U$, $\Sphere V$ of codimension $l\ge2$,
$k\ge2$ is connected (a path can be pushed off a submanifold of codimension at least two), so the
lift is connected. Hence there are exactly three boundary hypersurfaces, identified with those
of $\mathcal D_2$ as above: the lifts of $\Sphere V\times\RC U$ and $\RC V\times\Sphere U$ and the front
face over the corner $\Sphere V\times\Sphere U$. When $k=1$ or $l=1$ some of these split into two
components, as in the octagon.
